\documentclass[11pt,a4paper]{article}
\usepackage[T1]{fontenc}
\usepackage[utf8]{inputenc}
\usepackage{lmodern}
\usepackage[margin=27mm,headheight=15pt]{geometry}
\usepackage{amsmath,amssymb,amsthm,mathtools}
\usepackage{microtype,booktabs,array,longtable,enumitem,aliascnt}
\usepackage{xurl}
\usepackage[hidelinks,pdfusetitle]{hyperref}
\usepackage[nameinlink,noabbrev]{cleveref}
\hypersetup{pdftitle={Finite Residue Obstructions and Logarithmic-Depth Promotion in Hahn Transseries},pdfsubject={Explicit finite-dimensional obstructions and sharp logarithmic degree bounds},pdfkeywords={Hahn series, transseries, residues, differential equations, logarithmic depth}}
\usepackage{fancyhdr}
\pagestyle{fancy}
\fancyhf{}
\fancyhead[L]{\small Finite residue obstructions in Hahn transseries}
\fancyhead[R]{\small Research article}
\fancyfoot[C]{\thepage}
\setlength{\parskip}{4pt}
\setlength{\parindent}{1.2em}
\setlength{\emergencystretch}{2em}
\numberwithin{equation}{section}
\newtheorem{theorem}{Theorem}[section]
\newaliascnt{lemma}{theorem}
\newtheorem{lemma}[lemma]{Lemma}
\aliascntresetthe{lemma}
\crefname{lemma}{lemma}{lemmas}
\newaliascnt{proposition}{theorem}
\newtheorem{proposition}[proposition]{Proposition}
\aliascntresetthe{proposition}
\crefname{proposition}{proposition}{propositions}
\newaliascnt{corollary}{theorem}
\newtheorem{corollary}[corollary]{Corollary}
\aliascntresetthe{corollary}
\crefname{corollary}{corollary}{corollaries}
\theoremstyle{definition}
\newaliascnt{definition}{theorem}
\newtheorem{definition}[definition]{Definition}
\aliascntresetthe{definition}
\crefname{definition}{definition}{definitions}
\newaliascnt{example}{theorem}
\newtheorem{example}[example]{Example}
\aliascntresetthe{example}
\crefname{example}{example}{examples}
\theoremstyle{remark}
\newaliascnt{remark}{theorem}
\newtheorem{remark}[remark]{Remark}
\aliascntresetthe{remark}
\crefname{remark}{remark}{remarks}
% ed. (2026-09-29): unnumbered environment for ProveIt editorial notes; it
% leaves every numbered statement of the delivered article unchanged.
\newtheorem*{ednote}{Editorial note (ProveIt, 2026-09-29)}
\newcommand{\R}{\mathbb R}
\newcommand{\N}{\mathbb N}
\newcommand{\HH}{\mathcal H}
\newcommand{\cL}{\mathcal L}
\newcommand{\supp}{\operatorname{supp}}
\newcommand{\Res}{\operatorname{Res}}
\newcommand{\im}{\operatorname{im}}
\newcommand{\coker}{\operatorname{coker}}
\newcommand{\rank}{\operatorname{rank}}
\newcommand{\id}{\operatorname{id}}
\newcommand{\coef}[2]{[#1]#2}
\newcommand{\dprim}{d_0}
\newcommand{\eps}{\varepsilon}
\newcommand{\shift}{\mathcal N}
\newcommand{\snap}{\texttt{8813fb3a2cdd1702d3ea7f7ca5d63ae710630bef}}
\title{\textbf{Finite Residue Obstructions and\\Logarithmic-Depth Promotion\\in Hahn Transseries}\\[0.6em]
\large An explicit Fredholm reduction for small\\Euler-type differential perturbations}
\author{Research prepared for Vladimir Reshetnikov}
\date{29 September 2026}
\begin{document}
\maketitle
\thispagestyle{empty}
\begin{abstract}
At a fixed logarithmic depth, a formal differential equation can fail to have a solution even when all its nonresonant blocks are explicitly invertible. We give a concrete obstruction calculus for this failure. In the real Hahn field generated by $x,\log x,\ldots,\log_n x$, the image of the $m$th derivative is characterized by exactly $m$ residue moments. For an Euler-type operator $P(xD)$ perturbed by differential terms whose coefficients uniformly lower the exponent of $x$, all solvability conditions reduce to a finite real matrix. Its entries involve only finitely many perturbation steps, with an explicit cutoff determined by the real indicial roots and the support decrease. Its rank measures exactly the loss of homogeneous solutions and the remaining inhomogeneous obstructions at the original depth.

Adjoining one further logarithm removes every obstruction in this class. More precisely, if $P$ has $s$ distinct real roots, a particular solution has degree at most $s$ in the new logarithm, and every homogeneous solution has degree at most $s-1$. These bounds depend on the number of distinct roots, not their multiplicities. A family $\prod_{j=0}^{s-1}(xD-j)+(x\log x\cdots\log_n x)^{-1}$ attains both bounds and has a single Jordan-shift obstruction matrix. For its second-order member we construct actual real solutions and give an explicit error bound for the forced $\log\log x$ block, including a signed factorial-series remainder. The paper provides conventional proofs, finite exact verification scripts, a repository crosswalk, and a program of further questions.
\end{abstract}
\vfill
\noindent\fbox{\begin{minipage}{0.94\linewidth}\small
\textbf{Scope and status.} This is an unrefereed, self-contained research manuscript, not a Lean-verified development. Classical Hahn support calculus, block inversion, and perturbative elimination are acknowledged explicitly. The new contribution proposed here is their concrete residue-matrix and sharp logarithmic-degree synthesis for the stated operator class. No exhaustive priority claim, general transseries solvability theorem, or general analytic summability theorem is asserted.
\end{minipage}}
\clearpage
\setcounter{tocdepth}{1}
\tableofcontents
\vspace{1.2em}
\section*{Reading guide}
For the main classification theorem, start with \cref{thm:fredholm} and the
finite cutoff in \cref{thm:cutoff}. The logarithmic-depth and degree theorem
is \cref{thm:promotion}; its sharpness examples are
\cref{thm:sharp-matrix,thm:sharp-degree,cor:sharp-forcing}.

The logical route is: normalized integration and residue moments
(\cref{sec:residue,sec:moments}); the finite splitting of the Euler operator
(\cref{sec:euler}); perturbative elimination and support bounds
(\cref{sec:fredholm}); one-logarithm promotion and sharp examples
(\cref{sec:promotion,sec:sharp}). These formal arguments do not require an
analytic realization theorem.

The real-variable result can be read separately in \cref{sec:analytic}.
Its endpoint is the explicit inequality \eqref{eq:analytic-bound}, proved
by a Volterra integral argument. Verification scope and the proposed Lean
implementation are discussed in \cref{sec:verification}; ten further
research directions are developed in \cref{sec:questions}.
\clearpage

\section{Research target, repository context, and main conclusions}
\label{sec:introduction}

\subsection{The question being answered}
Formal antidifferentiation is often described by saying that one divides a monomial by its exponent, adding a logarithm at resonance. This description is useful but incomplete for a differential equation. Several resonant blocks may interact, the interaction may be mediated by infinitely many smaller terms, and it is not automatic that the resulting support is admissible. Nor does a formal solution by itself provide a real function with a certified error estimate.

This article isolates a class in which these questions have exact answers. The coefficients may be arbitrary real Hahn series of finite logarithmic depth; their supports need not be finitely generated. The essential restriction is that every perturbing coefficient decreases the \emph{outer exponent of $x$} by a common positive amount. Inside a fixed $x$-block, arbitrary real powers of the logarithms remain available.

The central questions are as follows. Which forcing terms can be solved without increasing logarithmic depth? Can infinitely supported perturbations be reduced to finite obstruction data? How many new logarithms, and how large a power of the new logarithm, are necessary? These are formulated and resolved below for a precise class, rather than presented as a solution of an unspecified broad conjecture.

\subsection{Relation to ProveIt}
The inspected ProveIt development separates the transseries library from the Fabius-function library. Its current transseries directory contains asymptotic scales, well-based supports, polynomial--logarithmic blocks, and inversion material \cite{repo-root,repo-volume}. The actual module
\texttt{TransseriesBlockAntiderivative.lean}
proves an explicit inverse for $\partial_L-c$ on polynomials when $c\ne0$, and proves polynomial antidifferentiation at $c=0$ \cite{repo-block}. Its finite inverse is
\begin{equation}
 (\partial_L-c)^{-1}f
 =-\sum_{k=0}^{\deg f}c^{-k-1}\partial_L^k f.
 \label{eq:repo-inverse}
\end{equation}
The documentation carefully distinguishes these results from the remaining construction of a concrete Laurent/Hahn ambient and associated uniqueness and evaluation statements \cite{repo-volume}.

We extend the mathematical setting from polynomial blocks to well-based Hahn blocks, replace the finite sum in \eqref{eq:repo-inverse} by a strongly summable operator series when needed, and retain the resonant defects instead of hiding them in an enlarged ambient field. This produces a finite matrix that decides whether the enlargement is actually necessary. We do not claim that the existing Lean results already prove the theorems of this article, and we have not modified or rebuilt the repository.

The provenance snapshot used for the repository references is
\begin{center}\small\snap.\end{center}
The transseries destination is \texttt{Analysis/Transseries/}, not the historical Fabius documentation directory. The investigation was targeted at the relevant documentation and modules; it was not an exhaustive audit of every source in the repository.

\subsection{Classical ingredients and the proposed contribution}
Hahn fields, strong series derivations, and general integration criteria have an established theory. In particular, Kuhlmann and Matusinski study strong derivations of Hardy type and give a general surjectivity criterion \cite{km2012}. The one-residue integration result below is an explicit finite-depth specialization of that surrounding theory, not a claim to have first discovered formal integration. Van der Hoeven's operator theory already supplies powerful implicit-function methods for generalized series \cite{vdh2001}, and his monograph develops a broad transseries differential algebra \cite{vdh2006}. Our geometric operator inverses are deliberately elementary instances of support-contracting methods.

Likewise, passing from an invertible complement to a finite effective operator is standard perturbative elimination; it is related in structure to homological perturbation methods \cite{hogancamp}. We prove the needed algebra directly. What is proposed here as a useful research extension is the \emph{explicit combination} of residue moments, a sharp finite perturbation cutoff, exact fixed-depth kernel/cokernel dimensions, a distinct-root bound for the power of one newly adjoined logarithm, and a family attaining the bounds. The literature check establishes relevant precedents, not priority for every formulation in this synthesis.

\subsection{Roadmap of the results}
Write $E=xD$ and let $P$ be a nonzero real polynomial. Let $r$ be the sum of the multiplicities of its real roots and $s$ their number. For
\begin{equation}
 \cL=P(E)+R,\qquad
 R=\sum_{j=0}^{\deg P}a_jE^j,
 \qquad \dprim(a_j)\le-\eps<0,
 \label{eq:main-class}
\end{equation}
the conclusions are:
\begin{enumerate}[label=(\roman*),leftmargin=2em,itemsep=3pt]
\item A canonically normalized block construction gives a finite matrix $M$ and a surjective defect map $b$ such that $\cL u=f$ at depth $n$ if and only if $Mc=b(f)$ is solvable. Both kernel and cokernel have dimension $r-\rank M$.
\item If $\Delta$ is the difference between the largest and smallest real roots, the matrix uses at most $\lfloor\Delta/\eps\rfloor$ perturbation factors. This is an exact cutoff, not an asymptotic approximation.
\item In the next logarithmic field there is a particular solution polynomial of degree at most $s$ in the new logarithm; all homogeneous solutions are polynomial of degree at most $s-1$. Explicit examples attain these bounds.
\end{enumerate}
The formal theorems are in \cref{sec:fredholm,sec:promotion,sec:sharp}. An actual differential equation with a quantitative real-variable certificate is treated in \cref{sec:analytic}.

\section{The finite-depth Hahn model}
\label{sec:model}

\subsection{Monomials, support, and summability}
For $n\ge0$, introduce formal symbols
\[
 \ell_0=x,\qquad \ell_1=\log x,\qquad\ldots,\qquad\ell_n=\log_n x.
\]
These names specify the intended differential structure, not an assertion that every infinite formal sum converges after real evaluation. For $a\in\R^{n+1}$ put
\[
 \ell^a=\prod_{j=0}^{n}\ell_j^{a_j}.
\]
Order exponents lexicographically, with $a_0$ the most significant coordinate. Larger exponent vectors mean larger monomials. Define
\begin{equation}
 \HH_n=\left\{\sum_{a\in S}c_a\ell^a:
 c_a\in\R,\ S\subset\R^{n+1}\text{ is reverse well-ordered}\right\}.
 \label{eq:hahn-field}
\end{equation}
Reverse well-ordered means that every nonempty subset has a greatest element. Coefficients may vanish; support always means the set of nonzero coefficients. We use the large-to-small convention throughout. Reversing the exponent order gives the usual valuation convention for Hahn series.

A family $(f_i)_{i\in I}$ is \emph{summable} if the union of its supports is reverse well-ordered and each monomial occurs in only finitely many members. Its sum is then defined coefficientwise. A linear operator is \emph{strong} if it preserves summable families and commutes with their sums.

We use the standard Hahn support lemma: finite unions and finite sums of reverse well-ordered subsets of an ordered abelian group are reverse well-ordered, and a fixed exponent has only finitely many decompositions as a sum of elements from two fixed such supports. The positive-support version of Neumann's lemma also makes geometric inverses of small Hahn series legitimate. These facts give the usual field operations in \eqref{eq:hahn-field}; they are part of the classical support calculus \cite{vdh2001,vdh2006}. The arguments below specify the additional support bounds needed for every operator series used in the paper.

\subsection{Outer blocks}
For $n\ge1$, let
\[
 L=\ell_1,\qquad
 \mathcal B_n=\R((L^\R\ell_2^\R\cdots\ell_n^\R)).
\]
Thus $\mathcal B_n$ is a copy of $\HH_{n-1}$ whose first variable is $L$. Each element of $\HH_n$ has a unique block decomposition
\begin{equation}
 f=\sum_{\alpha}x^\alpha f_\alpha,\qquad f_\alpha\in\mathcal B_n.
 \label{eq:blocks}
\end{equation}
The set of nonzero outer exponents is reverse well-ordered, and each block has reverse well-ordered lower support. Conversely these two properties imply reverse well-ordering of the full lexicographic support. Indeed, to find the greatest element of a nonempty subset, first choose its greatest outer exponent and then the greatest element of that fibre. The same argument proves the projection and fibre assertions in the other direction.

For $f\ne0$, define
\[
 \dprim(f)=\max\{\alpha:f_\alpha\ne0\},\qquad \dprim(0)=-\infty.
\]
Then $\dprim(fg)=\dprim(f)+\dprim(g)$ for nonzero products, and $\dprim(f+g)\le\max(\dprim(f),\dprim(g))$. These statements concern support only; they are not estimates for real-valued functions.

\subsection{The strong derivative}
Define on monomials
\begin{equation}
 D\ell^a=\sum_{j=0}^{n}a_j\ell^{a-\mathbf1_{\{0,\ldots,j\}}}.
 \label{eq:monomial-derivative}
\end{equation}
In particular, $D\ell_0=1$ and
\[
 D\ell_j=(\ell_0\ell_1\cdots\ell_{j-1})^{-1}\quad(j\ge1).
\]
There are only $n+1$ possible support shifts in \eqref{eq:monomial-derivative}. Their finite union is admissible, and a fixed output monomial receives at most $n+1$ input contributions. Consequently $D$ extends strongly to $\HH_n$. The Leibniz rule holds on monomials and then on arbitrary series by summability of Hahn products.

Let $\delta=d/dL$ denote the same derivative construction on $\mathcal B_n$. Directly from \eqref{eq:monomial-derivative},
\begin{align}
 D(x^\alpha A)&=x^{\alpha-1}(\alpha+\delta)A,\\
 E(x^\alpha A)&=x^\alpha(\alpha+\delta)A,
 \qquad E=xD.
 \label{eq:euler-block}
\end{align}
Every term of $\delta A$ has its $L$-exponent lowered by one. This elementary fact replaces analytic norm estimates inside a block.

\begin{lemma}[Formal functional calculus inside a block]
\label{lem:functional}
For $F(z)=\sum_{k\ge0}F_kz^k\in\R[[z]]$, the operator
$F(\delta)A=\sum_{k\ge0}F_k\delta^kA$ is well-defined and strong on $\mathcal B_n$. The substitution respects sums and products of formal power series. If $F(0)\ne0$, $F(\delta)$ is invertible with inverse $(1/F)(\delta)$.
\end{lemma}
\begin{proof}
If the largest $L$-exponent of $A$ is $B$, then that of $\delta^kA$ is at most $B-k$. At any fixed $L$-exponent, only finitely many $k$ can contribute. The union of supports is reverse well-ordered: above any fixed lower bound for the $L$-exponent only finitely many translated supports are involved. The same reasoning applies to a summable family because its union has a single upper bound $B$. It justifies rearranging the double sums in multiplication of formal operator series. Substituting the formal identity $F(z)(1/F(z))=1$ proves the inverse assertion.
\end{proof}

\begin{remark}
This argument requires no positive lower bound on $|F(0)|$ as parameters vary from one outer block to another. Small denominators may matter for analytic realization, but do not invalidate the formal block construction. An analytic estimate must therefore not be inferred from this lemma.
\end{remark}

\section{A residue exact sequence and normalized integration}
\label{sec:residue}
Put
\[
 \rho_n=(x\ell_1\cdots\ell_n)^{-1},\qquad
 \Res_n(f)=\coef{\rho_n}{f}.
\]
The coefficient of the constant monomial is denoted by $[1]f$.

\begin{theorem}[The fixed-depth integration defect]
\label{thm:residue}
There is a unique strong linear operator $I_n:\HH_n\to\HH_n$ satisfying
\begin{align}
 DI_nf&=f-\Res_n(f)\rho_n,&[1]I_nf&=0,\\
 I_nDf&=f-[1]f.&&
 \label{eq:integration-homotopy}
\end{align}
In particular,
\begin{equation}
 0\longrightarrow\R\longrightarrow\HH_n
 \xrightarrow{D}\HH_n\xrightarrow{\Res_n}\R
 \longrightarrow0
 \label{eq:exact-residue}
\end{equation}
is exact. A primitive of $f$ lies in $\HH_n$ exactly when $\Res_n(f)=0$. In $\HH_{n+1}$ a primitive is
\begin{equation}
 J_nf=I_nf+\Res_n(f)\ell_{n+1}.
 \label{eq:promoted-primitive}
\end{equation}
All primitives differ by a real constant.
\end{theorem}
\begin{proof}
For $n=0$, integrate $x^a$ as $x^{a+1}/(a+1)$ when $a\ne-1$, and send $x^{-1}$ to zero. The exceptional coefficient is exactly $\Res_0$. Translation of support proves strongness, and the asserted identities follow termwise.

Suppose the result has been proved for the lower field $\mathcal B_n$. In \eqref{eq:blocks}, write $I^{(L)}$ for its normalized integration operator. Define the primitive of a block with $\alpha\ne-1$ by
\begin{equation}
 I_n(x^\alpha A)=x^{\alpha+1}
 \sum_{k\ge0}\frac{(-1)^k}{(\alpha+1)^{k+1}}\delta^kA.
 \label{eq:nonresonant-integral}
\end{equation}
The sum is the inverse of $\alpha+1+\delta$ from \cref{lem:functional}. At $\alpha=-1$ define
\[
 I_n(x^{-1}A)=I^{(L)}A.
\]
The missing lower residue is the coefficient of
$\sigma_n=(L\ell_2\cdots\ell_n)^{-1}$ in $A$, and its outer counterpart is $\rho_n=x^{-1}\sigma_n$. This proves the first identity in \eqref{eq:integration-homotopy}.

Each input outer exponent is shifted to exactly one possible output outer exponent, namely $\alpha+1$. Within a nonresonant block the construction is strong by \cref{lem:functional}, and within the exceptional block it is strong by induction. The projection/fibre criterion therefore proves strongness on the full field. No constant is produced in a nonresonant block, and the exceptional block has been normalized to have zero constant.

To compute the kernel of $D$, apply \eqref{eq:euler-block} to a putative constant. A block with outer exponent $\alpha\ne0$ must vanish because $\alpha+\delta$ is invertible. The remaining outer-zero block lies in the kernel of $\delta$, which is $\R$ by induction. Thus $\ker D=\R$. Directly from the monomial derivative, $\Res_n(Df)=0$: every possible input that could produce $\rho_n$ has zero exponent in the coordinate being differentiated, and hence contributes zero. Hence $I_nDf$ and $f$ differ by a constant; normalization identifies it as $[1]f$.

The functional $\Res_n$ is onto because $\Res_n\rho_n=1$, proving exactness. Finally $D\ell_{n+1}=\rho_n$. The formula \eqref{eq:promoted-primitive} and the kernel calculation at depth $n+1$ prove the last assertions. Uniqueness of $I_n$ follows from its differential identity and zero-constant normalization.
\end{proof}

% ed. (2026-09-29): credits the canonical-volume theorem that this result
% re-derives and strengthens; the article does not cite it.
\begin{ednote}\label{ed:tower-strict}
Part of \cref{thm:residue} is already Theorem \path{plt:thm:ext-tower-strict} of the canonical volume (\path{Analysis/Transseries/docs/series-and-transseries/Transseries_And_Inversion/transseries_and_inversion.tex}), which this article does not cite. That theorem works in $\mathcal H_r(\mathbb K,\Gamma)$ with $\mathbb Z\subseteq\Gamma\subseteq\R$ and the monomial $\mathfrak n_r=L_0L_1\cdots L_r$, $L_0=x$, $L_{j+1}=\log L_j$, so that $\mathfrak n_n^{-1}=\rho_n$. Its equation \path{plt:eq:ext-residue}, $[\mathfrak n_r^{-1}]\partial_rF=0$ for all $F\in\mathcal H_r$, is the identity $\Res_n\circ D=0$, and its consequence that $L_{r+1}\notin\mathcal H_r$ is the statement here that $\ell_{n+1}$ is a primitive of $\rho_n$ outside $\HH_n$. \Cref{thm:residue} re-derives both for real coefficients and strengthens them to the exactness of \eqref{eq:exact-residue}: every $f$ with $\Res_n(f)=0$ has a primitive in $\HH_n$, given by the explicit operator $I_n$; \cref{thm:moments} extends the criterion to $D^m$. Neither result bears on the canonical volume's open Remark \path{plt:rmk:ext-open-depth}, which concerns compositional inverses rather than differential operators.
\end{ednote}

\begin{corollary}[Residue integration by parts]
\label{cor:ibp}
For $f,g\in\HH_n$,
\begin{equation}
 \Res_n(fDg)=-\Res_n(gDf).
 \label{eq:ibp}
\end{equation}
\end{corollary}
\begin{proof}
Apply $\Res_n$ to $D(fg)$ and use \cref{thm:residue}.
\end{proof}

\begin{example}
At depth one, $\rho_1=1/(x\log x)$. Thus this monomial has no primitive in $\HH_1$ and has primitive $\log\log x$ in $\HH_2$. In contrast, $x^{-1}(\log x)^{-2}$ integrates inside $\HH_1$ to $-(\log x)^{-1}$. The issue is a single exact coefficient, not merely the presence of a negative logarithmic power.
\end{example}

\section{Higher derivatives and residue moments}
\label{sec:moments}

\begin{theorem}[The complete obstruction to an $m$th primitive]
\label{thm:moments}
For every integer $m\ge1$,
\begin{equation}
 f\in D^m\HH_n
 \quad\Longleftrightarrow\quad
 \Res_n(x^k f)=0\quad(0\le k<m).
 \label{eq:moment-test}
\end{equation}
The kernel is $\R[x]_{<m}$, and the cokernel has dimension $m$ over $\R$. With $T=\ell_{n+1}$, a particular $m$th primitive is
\begin{equation}
 Y_m(f)=\frac1{(m-1)!}
 \sum_{k=0}^{m-1}(-1)^k\binom{m-1}{k}
 x^{m-1-k}J_n(x^k f).
 \label{eq:green}
\end{equation}
It has the form $U_m(f)+T V_m(f)$, with $U_m(f)\in\HH_n$ and
\begin{equation}
 V_m(f)=\frac1{(m-1)!}
 \sum_{k=0}^{m-1}(-1)^k\binom{m-1}{k}
 x^{m-1-k}\Res_n(x^k f).
 \label{eq:moment-polynomial}
\end{equation}
Thus an arbitrary number of repeated integrations requires at most the first power of one additional logarithm.
\end{theorem}
\begin{proof}
If $f=D^m u$, repeated use of \eqref{eq:ibp} gives
$\Res_n(x^kD^m u)=(-1)^m\Res_n(uD^m x^k)=0$ for $k<m$.

We verify the Green formula. For $m=1$, it is $Y_1=J_n$, and $DY_1(f)=f$. For $m>1$, differentiating a primitive factor in \eqref{eq:green} contributes
\[
 \frac{x^{m-1}f}{(m-1)!}
 \sum_{k=0}^{m-1}(-1)^k\binom{m-1}{k}=0.
\]
Differentiating the power of $x$ instead gives $Y_{m-1}(f)$, because
\[
 (m-1-k)\binom{m-1}{k}=(m-1)\binom{m-2}{k}.
\]
Therefore $D^mY_m(f)=f$. Substitution of \eqref{eq:promoted-primitive} yields \eqref{eq:moment-polynomial}. If the moments vanish, $Y_m(f)\in\HH_n$, proving sufficiency in \eqref{eq:moment-test}.

For the kernel, induct on $m$. If $D^m u=0$, then $D^{m-1}u$ is constant; subtract a suitable multiple of $x^{m-1}$ and use the induction hypothesis. Finally, the moment functionals are independent: their values on $x^{-j}\rho_n$, $0\le j<m$, form the identity matrix. This gives cokernel dimension $m$. The distinct powers of $x$ in \eqref{eq:moment-polynomial} also show that $V_m(f)=0$ exactly when all moments vanish.
\end{proof}

\begin{remark}[Why successive naive integrations obscure the answer]
One may integrate $m$ times and introduce a logarithm at each stage. Formula \eqref{eq:green} proves that this is unnecessary. The entire defect is already a finite polynomial of residue moments multiplying one new logarithm. This uniform control is the reason that repeated indicial roots will not increase the logarithmic-degree bound in \cref{sec:promotion}.
\end{remark}

\subsection{Integration of polynomials in the new logarithm}
We need a slightly stronger statement for later perturbative steps.

\begin{lemma}[One-step degree increase]
\label{lem:polynomial-extension}
On $\HH_n[T]\subset\HH_{n+1}$, with $T=\ell_{n+1}$, the derivative is surjective. It admits a linear right inverse that raises $T$-degree by at most one and is strong on each subspace of bounded $T$-degree. For every $m\ge1$, $D^m$ also admits such a right inverse, still raising $T$-degree by at most one, independently of $m$.
\end{lemma}
\begin{proof}
Let the highest-degree term of a polynomial forcing be $a_dT^d$. The derivative of
\[
 I_n(a_d)T^d+\frac{\Res_n(a_d)}{d+1}T^{d+1}
\]
is
\[
 a_dT^d+d\rho_nI_n(a_d)T^{d-1}.
\]
Subtract the second term from the remaining forcing and repeat at the next lower degree. The process ends after finitely many polynomial degrees, is linear, and uses only strong operations. Its degree increase is at most one. Apply the Green formula \eqref{eq:green}, now using this right inverse for each polynomial forcing $x^kf$, to obtain a right inverse of $D^m$ with the same degree bound. The verification of that formula used only $DJ=\id$, so applies unchanged.
\end{proof}

\section{The unperturbed Euler operator and its finite defect}
\label{sec:euler}
Throughout the rest of the formal analysis let $n\ge1$. Write the distinct real roots of a nonzero polynomial $P\in\R[z]$ as
\[
 \lambda_1<\cdots<\lambda_s,
\]
with multiplicities $m_1,\ldots,m_s$. Set $r=\sum_{i=1}^s m_i$; if there are no real roots, $s=r=0$. Nonreal roots are allowed. They do not generate real-power homogeneous monomials in the present field.

For a root $\lambda$, factor
\begin{equation}
 P(\lambda+z)=z^{m_\lambda}Q_\lambda(z),
 \qquad Q_\lambda(0)\ne0.
 \label{eq:root-factorization}
\end{equation}
On $x^\alpha\mathcal B_n$, the operator $P_0=P(E)$ acts as $P(\alpha+\delta)$.

\subsection{Kernel coordinates and residue coordinates}
Let $V=\R^r$, with coordinates indexed by pairs $(\lambda,k)$ where $0\le k<m_\lambda$. Define
\begin{align}
 Kc&=\sum_{\lambda}\sum_{k=0}^{m_\lambda-1}
 c_{\lambda,k}x^\lambda L^k,\\
 (\Lambda f)_{\lambda,k}
 &=\coef{x^\lambda L^k\ell_2^0\cdots\ell_n^0}{f},\\
 (\Pi f)_{\lambda,k}
 &=\Res^{(L)}(L^k f_\lambda),\\
 Je_{\lambda,k}&=x^\lambda L^{-k}\sigma_n,
 \qquad\sigma_n=(L\ell_2\cdots\ell_n)^{-1}.
 \label{eq:splitting-maps}
\end{align}
Here $\Res^{(L)}$ is the residue functional of the lower field. The maps $K,J$ have different purposes: $K$ represents homogeneous solutions, whereas $J$ represents obstructions. Direct coefficient extraction gives
\begin{equation}
 \Lambda K=\id_V,\qquad \Pi J=\id_V.
 \label{eq:sections}
\end{equation}

\begin{theorem}[Explicit splitting of $P(E)$]
\label{thm:euler-split}
The kernel of $P_0$ is $KV$, and its image is $\ker\Pi$. There is a unique normalized strong linear operator $G:\HH_n\to\HH_n$ such that
\begin{align}
 P_0G&=\id-J\Pi,&GP_0&=\id-K\Lambda,\\
 \Lambda G&=0,&GJ&=0.
 \label{eq:split-identities}
\end{align}
It preserves each outer $x$-exponent block. Thus both the kernel and cokernel of $P_0$ have dimension $r$.
\end{theorem}
\begin{proof}
If $P(\alpha)\ne0$, \cref{lem:functional} gives the inverse
\begin{equation}
 [P(\alpha+\delta)]^{-1}
 =\left.\frac1{P(\alpha+z)}\right|_{z=\delta}.
 \label{eq:nonresonant-euler}
\end{equation}
At a real root $\lambda$, \eqref{eq:root-factorization} gives the product $\delta^{m_\lambda}Q_\lambda(\delta)$. The second factor is invertible and commutes with $\delta$. Therefore it maps $\im\delta^{m_\lambda}$ onto itself: both it and its inverse commute with $\delta^{m_\lambda}$. By \cref{thm:moments}, the image is characterized by
\[
 \Res^{(L)}(L^k f_\lambda)=0\quad(0\le k<m_\lambda).
\]
The kernel is the space of polynomials in $L$ of degree less than $m_\lambda$. Indeed, this is the kernel of $\delta^{m_\lambda}$, and $Q_\lambda(\delta)$ is an invertible endomorphism of that finite polynomial space.

For arbitrary $f$, the series $f-J\Pi f$ satisfies all image conditions. Solve $P_0u=f-J\Pi f$ blockwise using \eqref{eq:nonresonant-euler}, and at a resonant block using $Q_\lambda(\delta)^{-1}$ and \eqref{eq:green}. Subtract $K\Lambda u$ to normalize. The result is the unique $Gf$ with $\Lambda Gf=0$. It is strong because the lower integration, residue, and functional-calculus operations are strong; there are only finitely many exceptional outer blocks.

This proves $P_0G=\id-J\Pi$. Applied to $P_0f$, the same normalized uniqueness gives $GP_0f=f-K\Lambda f$. If $f=Je$, then $f-J\Pi f=0$, so $GJe=0$. The construction never changes an outer exponent, proving the final support assertion.
\end{proof}

\begin{remark}
The residue coordinates in \eqref{eq:splitting-maps} do not need factors of $Q_\lambda(\delta)^{-1}$. Such factors would give an alternative basis of obstruction functionals, but the commutation argument in the proof shows that the unweighted moment test already characterizes the image. This simplifies the matrix in the next section.
\end{remark}

\section{Finite-dimensional reduction of a Hahn differential equation}
\label{sec:fredholm}

\subsection{The support-contracting hypothesis}
Let $d=\deg P$ and let
\begin{equation}
 R=\sum_{j=0}^{d}a_j E^j,\qquad a_j\in\HH_n.
 \label{eq:perturbation}
\end{equation}
Suppose that a fixed $\eps>0$ satisfies
\begin{equation}
 a_j\ne0\quad\Longrightarrow\quad\dprim(a_j)\le-\eps.
 \label{eq:strict-outer-drop}
\end{equation}
The derivatives act on the argument to their right, not on $a_j$. Since $E$ preserves an outer block, multiplication gives
\begin{equation}
 \dprim(Rf)\le\dprim(f)-\eps,
 \qquad \dprim(GRf)\le\dprim(f)-\eps.
 \label{eq:drop}
\end{equation}
No restriction is imposed on logarithmic exponents inside the coefficient blocks. Taking $a_d\ne0$ is permitted; the leading differential coefficient is still a nonzero perturbation of the leading constant coefficient of $P$.

\begin{lemma}[The geometric inverse]
\label{lem:geometric}
The operator
\begin{equation}
 \shift=(\id+GR)^{-1}
       =\sum_{q\ge0}(-GR)^q
 \label{eq:neumann}
\end{equation}
is well-defined and strong on $\HH_n$. It is a two-sided inverse.
\end{lemma}
\begin{proof}
For a fixed $f$ with $\dprim(f)=A$, the $q$th term has outer exponent at most $A-q\eps$. At any prescribed outer exponent only finitely many values of $q$ are possible. The union of the outer supports is reverse well-ordered because their maxima tend to $-\infty$, and each of its fibres is a finite union of admissible supports. The argument is uniform on a summable family, using the greatest outer exponent of its union. Thus the series is strong. The finite geometric identity has remainder $(-GR)^{N+1}f$; every coefficient of that remainder vanishes for sufficiently large $N$, proving both inverse identities.
\end{proof}

Define the \emph{resonance matrix} and the \emph{effective forcing} by
\begin{align}
 M&=\Pi R\shift K:V\longrightarrow V,\\
 b(f)&=\Pi f-\Pi R\shift Gf.
 \label{eq:effective-data}
\end{align}
All compositions in these formulas are now defined, including their infinite operator series.

\begin{theorem}[Finite-residue Fredholm alternative]
\label{thm:fredholm}
Under \eqref{eq:strict-outer-drop}, the equation
\begin{equation}
 (P(E)+R)u=f
 \label{eq:perturbed-equation}
\end{equation}
has a solution in $\HH_n$ if and only if
\begin{equation}
 Mc=b(f)
 \label{eq:matrix-equation}
\end{equation}
has a solution in $V=\R^r$. All solutions are given by
\begin{equation}
 u=\shift(Kc+Gf),\qquad Mc=b(f).
 \label{eq:reconstruction}
\end{equation}
The parameter is exactly $c=\Lambda u$. Moreover,
\begin{align}
 \ker(P(E)+R)&\cong\ker M,\\
 \coker(P(E)+R)&\cong V/\im M,\\
 \dim\ker(P(E)+R)&=\dim\coker(P(E)+R)=r-\rank M.
 \label{eq:index}
\end{align}
These are algebraic dimensions; no Banach-space Fredholm assertion is intended.
\end{theorem}
\begin{proof}
Apply $G$ to \eqref{eq:perturbed-equation} and use \eqref{eq:split-identities}:
\[
 u-K\Lambda u+GRu=Gf.
\]
With $c=\Lambda u$, the invertibility from \cref{lem:geometric} gives \eqref{eq:reconstruction}. Applying $\Pi$ to the original equation gives $\Pi Ru=\Pi f$, which after substitution is precisely \eqref{eq:matrix-equation}.

Conversely, suppose $Mc=b(f)$ and define $u$ by \eqref{eq:reconstruction}. Since $\Lambda G=0$, the equation $(\id+GR)u=Kc+Gf$ implies $\Lambda u=c$. The residual
$h=(P_0+R)u-f$ then satisfies $Gh=0$. The matrix equation says $\Pi h=0$. The identity $P_0G=\id-J\Pi$ now gives $h=P_0Gh+J\Pi h=0$.

For the cokernel, first note that $bJ=\id_V$, because $GJ=0$. Thus $b$ is surjective. For arbitrary $u$,
\[
 \shift G(P_0+R)u=u-\shift K\Lambda u,
\]
and hence
\[
 b((P_0+R)u)=\Pi R\shift K\Lambda u=M\Lambda u.
\]
The solvability equivalence already proved says that $b(f)\in\im M$ exactly when $f\in\im(P_0+R)$. Consequently $b$ induces the asserted cokernel isomorphism. Formula \eqref{eq:reconstruction}, together with $\Lambda u=c$, proves the kernel isomorphism. Rank--nullity in $V$ completes the proof.
\end{proof}

\subsection{An exact finite perturbation cutoff}
The solution reconstruction usually is an infinite Hahn series. The matrix itself is not an infinite perturbation problem.

\begin{theorem}[Finite cutoff and triangular structure]
\label{thm:cutoff}
Suppose $r>0$, and put
\[
 \lambda_{\min}=\lambda_1,\quad\lambda_{\max}=\lambda_s,
 \quad\Delta=\lambda_{\max}-\lambda_{\min},
 \quad Q=\left\lfloor\frac{\Delta}{\eps}\right\rfloor.
\]
Then
\begin{equation}
 M=\sum_{q=1}^{Q}(-1)^{q-1}\Pi R(GR)^{q-1}K.
 \label{eq:finite-M}
\end{equation}
An empty sum is zero. A matrix entry can be nonzero only from a source root $\lambda$ to a strictly smaller target root $\mu$. Thus $M$ is strictly block triangular and $M^s=0$.

For $f\ne0$ with $A=\dprim(f)$, the correction in $b(f)$ similarly terminates after at most
\begin{equation}
 Q_f=\max\left(0,\left\lfloor\frac{A-\lambda_{\min}}{\eps}\right\rfloor\right)
 \label{eq:forcing-cutoff}
\end{equation}
perturbation factors.
\end{theorem}
\begin{proof}
A term starting in the source block $x^\lambda$ and containing $q$ factors of $R$ can reach only outer exponents at most $\lambda-q\eps$. The maps $G$ preserve outer exponents, and $\Pi$ reads only the real root blocks. If $q\eps>\Delta$, no target root is reachable from any source root. This proves \eqref{eq:finite-M}. Every nonzero term also has target strictly below source, proving the triangular assertion. A product of $s$ such strict root decreases is impossible, so $M^s=0$.

For the forcing correction replace the source bound $\lambda_{\max}$ by $A$; $Gf$ has no greater outer exponent than $f$. This gives \eqref{eq:forcing-cutoff}.
\end{proof}

\begin{corollary}[A coefficient-support selection rule]
\label{cor:selection}
Let $S\subset(0,\infty)$ be the set of all outer drops $-\beta$ for which some coefficient $a_j$ has a nonzero $x^\beta$ block. A matrix block from $\lambda$ to $\mu$ vanishes unless
\begin{equation}
 \lambda-\mu\in S+S+\cdots+S
 \quad\text{for some positive finite number of summands}.
 \label{eq:selection}
\end{equation}
In particular, if no positive difference of real roots lies in this positive additive semigroup, then $M=0$.
\end{corollary}
\begin{proof}
Each perturbation factor contributes one outer drop from $S$, while all intervening operators preserve the outer exponent. The total drop along a contributing term must equal $\lambda-\mu$.
\end{proof}

\begin{remark}[Multiplicity and remaining dimension]
Strict block triangularity implies at least $\max_i m_i$ homogeneous solutions at the original depth whenever there are real roots. Indeed, at any chosen root level $i$, the columns up to that level have image contained in the strictly lower row blocks, while the remaining columns span a space of dimension at most the total multiplicity above that level. Thus $\rank M\le r-m_i$ for every $i$. If there is only one distinct real root, $M=0$, although not every forcing is solvable: the $r$ residue-moment obstructions remain.
\end{remark}

\subsection{A rigorous support truncation certificate}
\begin{proposition}
\label{prop:tail}
For $y=Kc+Gf\ne0$, let
\[
 u_N=\sum_{q=0}^{N}(-GR)^q y,
 \qquad u=\shift y.
\]
Then
\begin{equation}
 \dprim(u-u_N)\le\dprim(y)-(N+1)\eps.
 \label{eq:tail}
\end{equation}
Hence all coefficients with outer exponent greater than the right-hand side are exact after $N$ steps.
\end{proposition}
\begin{proof}
Every term in the omitted tail contains at least $N+1$ support decreases. Cancellation can only reduce, not increase, the greatest remaining outer exponent.
\end{proof}

The bound \eqref{eq:tail} is useful for symbolic computation, but it is not a numerical error bound. It controls the outer exponent only and does not bound the real sizes or growth of the coefficient blocks. The analytic certificate in \cref{sec:analytic} is obtained by a separate argument.

\subsection{What is and is not finite}
Equation \eqref{eq:finite-M} is finite in \emph{perturbation depth}. Individual coefficient blocks may still be infinitely supported, and arbitrary real Hahn coefficients need not be given by computable data. Therefore the theorem is not an unconditional finite-time algorithm on all set-theoretically possible inputs.

For an effective representation with exact coefficient extraction, exact outer-exponent arithmetic, and the requisite block operators, the formula becomes an algorithm. For finite rational supports, the included verification program illustrates part of this regime. For arbitrary noncomputable real supports, it remains a mathematical finite-dimensional characterization. This distinction is essential: finite-dimensional output does not by itself make its entries computable.

\section{One new logarithm suffices, with a sharp degree bound}
\label{sec:promotion}
Let $T=\ell_{n+1}$. The ring $\HH_n[T]$ is naturally a differential subring of $\HH_{n+1}$. Its elements have only finitely many nonnegative integer powers of $T$, although their coefficients can be arbitrary elements of $\HH_n$.

\begin{lemma}[A blockwise right inverse after promotion]
\label{lem:promoted-inverse}
There is a linear right inverse $\widetilde G$ of $P_0=P(E)$ on $\HH_n[T]$ with the following properties:
\begin{enumerate}[label=(\alph*),leftmargin=2em,itemsep=2pt]
\item it preserves the outer $x$-exponent;
\item at a nonresonant outer exponent it does not increase $T$-degree;
\item at a resonant outer exponent it raises $T$-degree by at most one, regardless of the root multiplicity;
\item it is strong on every family with uniformly bounded $T$-degree.
\end{enumerate}
It may be normalized by subtracting its kernel coordinates.
\end{lemma}
\begin{proof}
On the lower polynomial ring $\mathcal B_n[T]$, the derivative still lowers the $L$-exponent by one, including the term $\delta T=\sigma_n$. Thus the functional calculus of \cref{lem:functional} remains valid on each bounded-degree polynomial space. Differentiation never increases $T$-degree.

At a nonresonant exponent use \eqref{eq:nonresonant-euler}. At a root $\lambda$, use \eqref{eq:root-factorization}, first a right inverse of $\delta^{m_\lambda}$ from \cref{lem:polynomial-extension} in the lower variable $L$, and then $Q_\lambda(\delta)^{-1}$. This raises degree by at most one, not by $m_\lambda$. Strongness and preservation of the outer block follow from the construction. Subtracting the finitely many coordinates in the kernel does not affect any of these bounds.
\end{proof}

\begin{theorem}[Logarithmic-depth promotion]
\label{thm:promotion}
Let $\cL=P(E)+R$ satisfy \eqref{eq:strict-outer-drop}, and let $s$ and $r$ be as in \cref{sec:euler}.

For every $f\in\HH_n$, there is a solution of $\cL u=f$ in $\HH_n[T]$ with
\begin{equation}
 \deg_T u\le s.
 \label{eq:inhom-degree}
\end{equation}
The homogeneous solution space in $\HH_{n+1}$ has dimension $r$, and all its elements lie in $\HH_n[T]$. If $s\ge1$, they satisfy
\begin{equation}
 \deg_T u\le s-1.
 \label{eq:hom-degree}
\end{equation}
If $s=0$, the homogeneous kernel is zero and every forcing has a unique solution already in $\HH_n$.

Consequently the matrix test of \cref{thm:fredholm} gives the exact answer to whether depth $n$ suffices. Whenever it fails, exactly one further logarithmic level suffices. For homogeneous solutions, the number of dimensions missing at depth $n$ is precisely $\rank M$.
\end{theorem}
\begin{proof}
Define successively the polynomial terms
\begin{equation}
 u=\sum_{q\ge0}(-\widetilde G R)^q\widetilde G f.
 \label{eq:promoted-particular}
\end{equation}
We first prove a uniform degree bound, before summing in the polynomial ring. Every application of $R$ lowers the outer exponent by at least $\eps$. The operators $\widetilde G$ preserve it. Along any contributing sequence of blocks, the outer exponents at successive applications of $\widetilde G$ therefore strictly decrease. Such a sequence can visit each real root at most once. By \cref{lem:promoted-inverse}, only a visit to a root can increase the $T$-degree, and each visit increases it by at most one. Starting from $f$ of $T$-degree zero, all terms in \eqref{eq:promoted-particular} have degree at most $s$.

This path assertion is coefficientwise legitimate. At a fixed perturbation length, the Hahn product and strong block operators give summable contributions. At a fixed final outer exponent, only finitely many perturbation lengths can occur, by the same outer bound as in \cref{lem:geometric}. Thus the sum is well-defined in $\HH_{n+1}$, and its uniform degree bound puts it in $\HH_n[T]$. It satisfies
$u+\widetilde G Ru=\widetilde G f$; applying $P_0$ proves $\cL u=f$.

For a kernel seed $Kc$, set
\begin{equation}
 u_c=\sum_{q\ge0}(-\widetilde G R)^qKc.
 \label{eq:promoted-homogeneous}
\end{equation}
Applying $P_0$ gives $\cL u_c=0$. A seed in the root block $\lambda$ begins with $T$-degree zero and undergoes a perturbation before any use of $\widetilde G$. It can therefore encounter only roots strictly smaller than $\lambda$. There are at most $s-1$ such roots, proving \eqref{eq:hom-degree}. The leading outer block of the solution from $x^\lambda L^k$ is exactly that seed, so the $r$ resulting solutions are linearly independent.

It remains to exclude additional homogeneous solutions in the full field $\HH_{n+1}$. Let $v\ne0$ solve $\cL v=0$, and let $\alpha$ be its greatest outer exponent. Since $Rv$ is strictly lower, the leading block $v_\alpha$ satisfies $P(\alpha+\delta)v_\alpha=0$ in the lower field with one more logarithmic variable. The same kernel proof as in \cref{thm:euler-split} shows that $\alpha$ is a real root and $v_\alpha$ is a polynomial in $L$ of degree less than its multiplicity. Subtract the corresponding linear combination of the solutions \eqref{eq:promoted-homogeneous}. The greatest possible remaining outer exponent is a strictly lower real root. Repeating removes the solution after at most $s$ root levels. Thus the constructed $r$ solutions span the entire kernel in $\HH_{n+1}$.

When there are no real roots, $G$ is an actual inverse of $P_0$, so \cref{thm:fredholm} gives the stated unique fixed-depth solution. Finally, a failed matrix test means no solution belongs to $\HH_n$, whereas \eqref{eq:promoted-particular} supplies one in the next field. The dimension statement follows by subtracting \eqref{eq:index} from $r$.
\end{proof}

\subsection{Distinct roots, not multiplicity}
The distinction in \eqref{eq:inhom-degree}--\eqref{eq:hom-degree} is significant. For example, when $P(z)=(z-\lambda)^m$, there is only one root level, however large $m$ is. All $m$ homogeneous solutions remain at the original logarithmic depth after an outer-small perturbation. A forcing may still need the new logarithm, but at most its first power. Repeated differentiation produces multiple residue moments, while increased degree in the \emph{new} logarithm is caused by traversing different resonant outer exponents.

This also explains why the two statements ``one new logarithm suffices'' and ``the correction is linear in that logarithm'' are not equivalent. The former is true throughout our class; the latter is false as soon as a long chain of distinct resonances is present.

\section{A sharp family with an exact Jordan-shift obstruction}
\label{sec:sharp}
For $n\ge1$ and an integer $s\ge1$, set
\begin{equation}
 P_s(z)=\prod_{j=0}^{s-1}(z-j),\qquad
 \cL_{n,s}=P_s(E)+\rho_n.
 \label{eq:sharp-family}
\end{equation}
Here the last term is multiplication by $\rho_n=x^{-1}\sigma_n$. The uniform support drop is $\eps=1$, and all roots are simple.

\begin{theorem}[Maximal loss of fixed-depth homogeneous solutions]
\label{thm:sharp-matrix}
In the bases $K e_j=x^j$ and $J e_j=x^j\sigma_n$, $0\le j<s$, the matrix of \cref{thm:fredholm} is exactly
\begin{equation}
 Me_0=0,\qquad Me_j=e_{j-1}\quad(1\le j<s).
 \label{eq:jordan}
\end{equation}
Consequently
\[
 \rank M=s-1,\qquad
 \dim\ker_{\HH_n}\cL_{n,s}
 =\dim\coker_{\HH_n}\cL_{n,s}=1,
\]
whereas $\dim\ker_{\HH_{n+1}}\cL_{n,s}=s$.
\end{theorem}
\begin{proof}
For $j\ge1$, multiplication by $\rho_n$ gives $RK e_j=J e_{j-1}$. Since $GJ=0$, every higher perturbation correction starting from that vector vanishes before the next $R$. The first term gives $\Pi RK e_j=e_{j-1}$. Starting from $e_0$, the first perturbation has outer exponent $-1$, and all later terms have still smaller outer exponent. No root block can be reached, so the column is zero. This proves \eqref{eq:jordan}. The dimension conclusions follow from \cref{thm:fredholm,thm:promotion}.
\end{proof}

\begin{theorem}[Sharp homogeneous logarithmic degree]
\label{thm:sharp-degree}
There is a homogeneous solution of \eqref{eq:sharp-family} whose leading outer block is $x^{s-1}$. In its block of outer exponent $s-1-k$, $0\le k<s$, the coefficient of the highest power $T^k$ is the nonzero real number
\begin{equation}
 c_0=1,\qquad
 c_k=-\frac{c_{k-1}}{kP_s'(s-1-k)}
 =\frac{(-1)^k}{k!\prod_{h=1}^k P_s'(s-1-h)}.
 \label{eq:sharp-coefficients}
\end{equation}
In particular, the bound $s-1$ in \eqref{eq:hom-degree} is attained. No homogeneous solution with that leading block can have smaller $T$-degree.
\end{theorem}
\begin{proof}
Use \eqref{eq:promoted-homogeneous} with seed $x^{s-1}$. Write the successive root blocks as $x^{s-1-k}A_k$. Since $R$ shifts the outer exponent by exactly $-1$, the root-block recursion is
\begin{equation}
 P_s(s-1-k+\delta)A_k=-\sigma_n A_{k-1},\qquad A_0=1.
 \label{eq:sharp-recursion}
\end{equation}
At $k\ge1$, factor the left operator as $\delta Q(\delta)$ with $Q(0)=P_s'(s-1-k)$. Suppose $A_{k-1}$ has top $T$-coefficient $c_{k-1}$, a real constant. The top contribution to a primitive of $-c_{k-1}\sigma_nT^{k-1}$ is $-c_{k-1}T^k/k$, since $\delta T=\sigma_n$. Terms of smaller $T$-degree can produce degree at most $k-1$.

Applying $Q(\delta)^{-1}$ divides the coefficient of $T^k$ by $Q(0)$. Every positive power of $\delta$ differentiates the constant coefficient times $T^k$ and lowers its $T$-degree. This proves the recurrence and shows that the top coefficient remains a real constant. It is nonzero because the roots are simple. Below the last root, all blocks are nonresonant, so no further degree increase occurs.

Two solutions with the same leading block differ by homogeneous solutions seeded at smaller roots. Such a solution can traverse at most $s-2$ lower resonances, so cannot cancel the nonzero $T^{s-1}$ coefficient at the last root. This proves the last assertion.
\end{proof}

Using
\begin{equation}
 P_s'(j)=(-1)^{s-1-j}j!(s-1-j)!,
 \label{eq:ps-derivative}
\end{equation}
the last coefficients for the first six cases are
\begin{center}
\begin{tabular}{ccl}
\toprule
$s$ & Necessary homogeneous degree & $c_{s-1}$ in the $x^0$ block\\
\midrule
1 & 0 & $1$\\
2 & 1 & $1$\\
3 & 2 & $-1/4$\\
4 & 3 & $-1/144$\\
5 & 4 & $1/82944$\\
6 & 5 & $1/1194393600$\\
\bottomrule
\end{tabular}
\end{center}
These rational identities were also checked by the included script, independently of the proof by recurrence.

\begin{corollary}[Sharp inhomogeneous degree]
\label{cor:sharp-forcing}
For the forcing
\begin{equation}
 f=x^{s-1}\sigma_n,
 \label{eq:sharp-forcing}
\end{equation}
every solution of $\cL_{n,s}u=f$ in $\HH_{n+1}$ has $T$-degree exactly $s$. The coefficient of $T^s$ in the outer-zero block is
\begin{equation}
 C_s=\frac{(-1)^{s-1+s(s-1)/2}}
 {s!\left(\prod_{j=0}^{s-1}j!\right)^2}.
 \label{eq:forced-coefficient}
\end{equation}
There is no solution in $\HH_n$.
\end{corollary}
\begin{proof}
At the highest root, the forcing produces the top term $T/P_s'(s-1)$. The same recursion as in \eqref{eq:sharp-recursion} then increases degree at each of the remaining $s-1$ roots. The top coefficient at the last root is
\[
 \frac{(-1)^{s-1}}{s!\prod_{j=0}^{s-1}P_s'(j)},
\]
which equals \eqref{eq:forced-coefficient} by \eqref{eq:ps-derivative}. Every homogeneous solution in the entire next-depth field has degree at most $s-1$ by \cref{thm:promotion}, so the coefficient cannot be canceled by choosing another particular solution.

The original-depth obstruction can also be read directly from the matrix. Here $f=Je_{s-1}$, so $b(f)=e_{s-1}$, which is not in the range of the Jordan shift \eqref{eq:jordan}.
\end{proof}

\subsection{Consequences for proposed closure principles}
The family gives a counterexample to the statement that outer-small perturbations of an Euler equation preserve all its homogeneous solutions in a fixed logarithmic field. The loss can be $s-1$ out of $s$ dimensions, even though the perturbation is only multiplication by one monomial. It also refutes a uniform bound on the degree of the newly needed logarithm that is independent of the number of distinct resonances.

By contrast, it does \emph{not} refute closure after one logarithmic-depth promotion: that is exactly what \cref{thm:promotion} proves. The distinction between depth and degree is the structural point of these examples.

\section{A second-order equation: formal resonance and real error bounds}
\label{sec:analytic}
We now pass from a formal theorem to an actual real-variable calculation. This section realizes the decisive first logarithmic block of one example, not every arbitrary Hahn series in the preceding sections.

Take $n=1$ and $s=2$. With $L=\log x$ and $T=\log L$, the equation is
\begin{equation}
 \biggl(E(E-1)+\frac1{xL}\biggr)y=0,
 \qquad\text{equivalently}\qquad
 y''+\frac{y}{x^3\log x}=0.
 \label{eq:analytic-ode}
\end{equation}
Its fixed-depth matrix is $Me_1=e_0$, $Me_0=0$. Exactly one homogeneous dimension remains in $\HH_1$, and the growing solution seeded by $x$ requires $T$.

\subsection{The exact formal outer-zero block}
The next block after $x$ solves
\begin{equation}
 \delta(\delta-1)A_0=-\frac1L.
 \label{eq:a0-equation}
\end{equation}
Define the formal factorial series
\begin{equation}
 \widehat S(L)=\sum_{k=0}^{\infty}(-1)^k k!L^{-k-1}.
 \label{eq:Sformal}
\end{equation}
It is a legitimate Hahn series, though divergent as a numerical power series for every fixed positive $L$. Termwise differentiation gives
\begin{equation}
 (\delta-1)\widehat S=-\frac1L.
 \label{eq:Sformal-ode}
\end{equation}
Since $\delta T=1/L$, one obtains
\[
 \delta(\delta-1)(T+\widehat S)=-\frac1L.
\]
The zero-constant choice for the outer-zero block is therefore
\begin{equation}
 y_{\mathrm{grow}}
 =x+T+\widehat S(L)+\text{blocks with outer exponent at most }-1.
 \label{eq:formal-growing}
\end{equation}
All lower blocks are polynomial of degree at most one in $T$ by \cref{thm:promotion}.

For comparison, the solution seeded by $1$ has the first correction
\begin{align}
 y_{\mathrm{bd}}&=1+x^{-1}v(L)+\text{lower outer blocks},\\
 v(L)&=-\sum_{k\ge0}(-1)^k k!\bigl(1-2^{-k-1}\bigr)L^{-k-1}\nonumber\\
 &=-\frac1{2L}+\frac3{4L^2}-\frac7{4L^3}+\cdots.
 \label{eq:bounded-formal}
\end{align}
Indeed, the block equation is $(\delta-1)(\delta-2)v=-1/L$, and
\[
 \frac1{(z-1)(z-2)}
 =\sum_{k\ge0}\bigl(1-2^{-k-1}\bigr)z^k.
\]
Neither \eqref{eq:formal-growing} nor \eqref{eq:bounded-formal} is yet a real convergence assertion.

\subsection{An exact sum and a signed factorial remainder}
For $L>0$, define
\begin{equation}
 S(L)=\int_0^\infty\frac{e^{-Lt}}{1+t}\,dt.
 \label{eq:laplace}
\end{equation}
Differentiation under the integral, justified by exponential domination on compact positive $L$-intervals, gives
\begin{equation}
 S'(L)-S(L)=-\frac1L,
 \qquad 0<S(L)<\frac1L.
 \label{eq:S-ode}
\end{equation}
Thus $S$ is an exact real counterpart of the formal block $\widehat S$.

\begin{proposition}[Signed remainder]
\label{prop:factorial-remainder}
For every integer $N\ge0$ and every $L>0$,
\begin{equation}
 0<(-1)^N\left(S(L)-\sum_{k=0}^{N-1}(-1)^kk!L^{-k-1}\right)
 <N!L^{-N-1}.
 \label{eq:factorial-remainder}
\end{equation}
The empty sum for $N=0$ is zero.
\end{proposition}
\begin{proof}
Use the finite geometric identity
\[
 \frac1{1+t}=\sum_{k=0}^{N-1}(-1)^kt^k+
 \frac{(-1)^Nt^N}{1+t}
\]
in \eqref{eq:laplace}. Each polynomial integral is $k!L^{-k-1}$. The remaining positive integral is strictly between zero and
$\int_0^\infty e^{-Lt}t^N\,dt=N!L^{-N-1}$.
\end{proof}
The Borel transform in the convention
$\sum a_kL^{-k-1}\mapsto\sum a_kt^k/k!$
is $1/(1+t)$. Formula \eqref{eq:laplace} is its Laplace integral on the positive axis. This establishes ordinary Borel summability of this single block in that direction. It does not establish Borel summability of the entire Hahn solution or arbitrary coefficient fields.

\subsection{Volterra construction of two real solutions}
Let $X\ge e^2$, and on $[X,\infty)$ put
\[
 q(x)=\frac1{x^3\log x},\qquad
 (Vg)(x)=\int_x^\infty(t-x)q(t)g(t)\,dt.
\]
For bounded continuous $g$,
\begin{equation}
 \|Vg\|_{[x,\infty)}\le\eta_x\|g\|_{[x,\infty)},
 \qquad \eta_x=\frac1{2x\log x}<1.
 \label{eq:Vbound}
\end{equation}
Here the norm is the supremum on the indicated tail. To prove this, use $\log t\ge\log x$ and
\[
 \int_x^\infty\frac{t-x}{t^3}\,dt=\frac1{2x}.
\]
The Neumann series for $(\id+V)^{-1}$ consequently converges in the usual supremum norm. Also $(Vg)''=qg$ for continuous bounded $g$, by differentiating the integral twice.

The bounded solution
\begin{equation}
 y_0=(\id+V)^{-1}1
 \label{eq:y0}
\end{equation}
satisfies $y_0''=-qy_0$ and
\begin{equation}
 |y_0(x)-1|\le\frac{\eta_x}{1-\eta_x}.
 \label{eq:y0-error}
\end{equation}
Thus $y_0(x)\to1$. The tail estimate follows by restricting the same Volterra equation to $[x,\infty)$ and applying \eqref{eq:Vbound}.

For the growing solution, define
\begin{equation}
 z(x)=x+\log\log x+S(\log x)=x+T+S(L).
 \label{eq:z}
\end{equation}
A direct calculation from \eqref{eq:S-ode} gives
\begin{equation}
 z''+qz=q(T+S(L)).
 \label{eq:z-residual}
\end{equation}
The function $T+S(L)$ is unbounded but its Volterra transform
$h=V(T+S(L))$ is bounded. Set
\begin{equation}
 w=-(\id+V)^{-1}h,\qquad y_1=z+w.
 \label{eq:y1}
\end{equation}
Then $w=-V(T+S(L))-Vw$, so $w''=-q(T+S(L)+w)$. Combining with \eqref{eq:z-residual} proves that $y_1$ solves \eqref{eq:analytic-ode} exactly.

\begin{theorem}[An explicit analytic certificate]
\label{thm:analytic}
For every $x\ge e^2$, the real solution $y_1$ constructed in \eqref{eq:y1} satisfies
\begin{equation}
 \left|y_1(x)-x-\log\log x-S(\log x)\right|
 \le
 \frac{\displaystyle\frac{\log\log x}{2x\log x}
       +\displaystyle\frac5{4x(\log x)^2}}
 {\displaystyle1-\frac1{2x\log x}}.
 \label{eq:analytic-bound}
\end{equation}
The solutions $y_0,y_1$ are linearly independent. For every fixed $N\ge0$,
\begin{align}
 y_1(x)={}&x+\log\log x+
 \sum_{k=0}^{N-1}(-1)^kk!(\log x)^{-k-1}
 +\mathcal E_N(x),\\
 |\mathcal E_N(x)|\le{}&
 N!(\log x)^{-N-1}
 +\frac{\displaystyle\frac{\log\log x}{2x\log x}
       +\displaystyle\frac5{4x(\log x)^2}}
 {\displaystyle1-\frac1{2x\log x}}.
 \label{eq:combined-bound}
\end{align}
\end{theorem}
\begin{proof}
The integrand defining $h$ is positive for $x\ge e^2$, and $h$ is decreasing. Substitute $t=xv$, with $L=\log x$ and $T=\log L$. For $v\ge1$,
\[
 \log\log(xv)=\log(L+\log v)\le T+\frac{\log v}{L},
 \qquad S(\log(xv))\le\frac1L.
\]
Consequently
\begin{align*}
 h(x)&\le\frac1{xL}\int_1^\infty
 \frac{v-1}{v^3}\left(T+\frac{\log v+1}{L}\right)dv\\
 &=\frac{T}{2xL}+\frac5{4xL^2},
\end{align*}
where
\[
 \int_1^\infty\frac{v-1}{v^3}\,dv=\frac12,
 \qquad
 \int_1^\infty\frac{v-1}{v^3}\log v\,dv=\frac34.
\]
Because $h$ is decreasing, its norm on $[x,\infty)$ is $h(x)$. The inverse bound from \eqref{eq:Vbound} therefore gives $|w(x)|\le h(x)/(1-\eta_x)$, proving \eqref{eq:analytic-bound}.

We have $y_0(x)\to1$, while $y_1(x)/x\to1$, so the solutions cannot be scalar multiples. The integral equations make them $C^2$ solutions, and the ordinary uniqueness theorem for a second-order linear equation shows that they constitute a basis on the interval. Finally combine \eqref{eq:analytic-bound} with \cref{prop:factorial-remainder} to obtain \eqref{eq:combined-bound}.
\end{proof}

The normalization implicit here is
\[
 y_0(x)\to1,\qquad y_1(x)-x-\log\log x\to0.
\]
Adding a multiple of $y_0$ to $y_1$ would change the second limit. This selects a definite real growing solution, rather than merely its leading asymptotic equivalence.

\subsection{Numerical illustrations and their interpretation}
The accompanying script evaluates $S$, $z$, and the right-hand side of \eqref{eq:analytic-bound} at high working precision. It records the following approximate values:
\begin{center}
\begin{tabular}{rccc}
\toprule
$x$ & $S(\log x)$ & $z(x)$ & Bound in \eqref{eq:analytic-bound}\\
\midrule
$10$ & $0.323897895933$ & $11.1579303412$ & $4.26125667\cdot10^{-2}$\\
$100$ & $0.182974349963$ & $101.710153976$ & $2.24996864\cdot10^{-3}$\\
$10^4$ & $0.0987515406201$ & $10002.3190783$ & $1.35270483\cdot10^{-5}$\\
\bottomrule
\end{tabular}
\end{center}
The table entries are rounded illustrations, not outward-rounded interval certificates. The exact formula \eqref{eq:analytic-bound}, proved above, is the certificate. No numerical integration of the differential equation is needed to justify it.

For fixed $N$, the second term on the right of \eqref{eq:combined-bound} is smaller than every fixed inverse power of $\log x$. This realizes the whole displayed factorial expansion within the decisive outer-zero block. We do not claim here that the construction has been matched term by term to every smaller outer block of the full formal solution.

\section{Proof boundaries, checks, and formalization}
\label{sec:verification}

\subsection{What the proofs establish}
The formal statements are algebraic statements about explicitly defined Hahn fields and strong operators. Their infinite sums are justified by two distinct mechanisms: inside a block, $\delta$ lowers the $L$-exponent by one; between blocks, $R$ lowers the $x$-exponent by a uniform $\eps$. Residue extraction accounts for the finite exceptional part. Neither mechanism uses a hidden analytic convergence assumption.

The analytic theorem concerns a concrete real equation and is established separately by a supremum-norm contraction. It does not follow merely by substituting a real number into an arbitrary formal Hahn series. The factorial block is divergent as an ordinary power series, but its explicit Laplace integral has a signed remainder bound.

No claim of Lean compilation, mechanized checking of the new theorems, or proof-assistant coverage is made. The repository modules motivated and supplied the mathematical starting point; the proofs in this article stand as conventional proofs requiring normal mathematical review.

\subsection{Finite exact verification supplied with the article}
The package contains \path{verification/verify.py} and its recorded output \path{verification/results.json}. The run completed with all checks passing. Its scope is intentionally narrower than theorem verification.

The exact symbolic part uses rational arithmetic to check 212 block identities. A finite block is represented as a dictionary of coefficients of $L^bT^c$, with integer $b$ and nonnegative integer $c$. Differentiation uses
\[
 \delta(L^bT^c)=bL^{b-1}T^c+cL^{b-1}T^{c-1}.
\]
The program tests resonant and nonresonant inverses, repeated-root cases, and the sharp-family recursion for $1\le s\le6$. It retains guard terms down to $L$-exponent $-28$ and checks residuals only at exponents at least $-10$. The gap prevents a discarded lower term from being mistaken for a checked identity at the visible cutoff.

A second exact test constructs 24 finite-dimensional split-operator examples, with deterministic random seed $20260929$. It checks the geometric inverse, finite matrix, reconstruction formula, and equality of kernel/cokernel dimensions using exact SymPy matrices. These examples validate the algebraic arrangement and signs; they do not replace the infinite-support proof.

The numerical part uses \texttt{mpmath} for the integral block and remainder illustrations. It is not an interval-arithmetic verification. Its values are reproducible at the script's chosen working precision. The symbolic dependency versions used in the delivered run are recorded in the package README.

\subsection{A staged Lean implementation plan}
The following separation would keep a future formalization auditable.

\paragraph{Stage 1: the concrete ambient and strong derivative.}
Construct the ordered exponent group $\R^{n+1}$ with lexicographic order and its order-dual support convention. Establish block decomposition, coefficient extraction, and the finite-shift derivative. This stage needs a concrete Hahn ambient, not only polynomial identities. The repository's well-based support lemmas and its existing polynomial block API are relevant starting points \cite{repo-volume,repo-block}.

\paragraph{Stage 2: normalized integration and residue moments.}
Formalize the lower-exponent contracting functional calculus, define $I_n$, and prove \eqref{eq:integration-homotopy}. Then derive residue integration by parts and the Green formula. The moment theorem should be proved abstractly for any differential algebra with $Dx=1$, constants $\R$, and the stated residue splitting; the Hahn construction supplies an instance.

\paragraph{Stage 3: a generic finite-defect perturbation theorem.}
Separate \cref{thm:fredholm} from the concrete field. Its proof only uses the identities \eqref{eq:split-identities}, a geometric inverse, and finite-dimensional linear algebra. This reusable algebraic theorem can be formalized before all transseries details are available. Its hypotheses must explicitly include the sum/existence conditions rather than assume them implicitly.

\paragraph{Stage 4: support bounds and logarithmic promotion.}
Instantiate the abstract theorem using the primary-exponent decrease. Prove the exact cutoff, the degree-controlled right inverse on $\HH_n[T]$, and the strictly descending root-path argument. This is the part that turns a general perturbative method into a sharp fixed-depth theorem.

\paragraph{Stage 5: examples and analytic certification.}
The Jordan-shift family is a compact integration test of the formal API. The Volterra theorem belongs in a separate real-analysis development with explicit differentiation-under-the-integral and norm bounds. Formal symbolic truncation and analytic error bounds should remain different theorem types.

This plan does not presume that a full class of real germs realizing arbitrary Hahn series can be formalized as a consequence. Such an evaluation theory would require additional hypotheses and constructions.

\section{Further research questions and concrete next targets}
\label{sec:questions}
The following are proposed questions arising from the proved theorems. They are not presented as a catalog of already established open problems in the literature. Each indicates what the present argument does not settle.

\subsection{The exact minimal power of the promoted logarithm}
The finite matrix decides whether degree zero in $T$ is possible; the root-path argument gives the universal upper bound $s$. It does not in general compute the least possible positive degree. Can one construct a finite sequence of residue matrices whose successive ranks recover the full filtration by $T$-degree of homogeneous and particular solutions?

A concrete first target is the subclass with finitely many rational outer drops and finite Laurent-log coefficient blocks. One would seek an algorithm that outputs the smallest $d$ such that a forcing has a solution in $\HH_n[T]_{\le d}$, together with an obstruction certificate for degree $d-1$. The sharp family should yield the complete sequence $0,1,\ldots,s$.

% ed. (2026-09-29): records the later package that answers the homogeneous
% half of this question in a different (depth-0 systems) setting.
\begin{ednote}\label{ed:filtration}
The later package \path{Analysis/Transseries/docs/series-and-transseries/Hahn_Fuchsian_Resonance_Analytic_Normalization/} answers the homogeneous half of this question in a neighbouring setting, without citing it (it knew this article only through its README). For first-order systems $xY'=(A_0+A_+(x))Y$ over the log-free Hahn field with $L=\log x$ adjoined, with constant coefficient matrices, real spectrum and well-ordered positive exponents, the solutions of degree at most $d$ in $L$ form a space of dimension $\dim\ker B^{d+1}$ for one finite nilpotent matrix $B=N+\sum_\rho R_\rho$ (\texttt{hahn\_fuchsian.tex}, Theorem \texttt{thm:logs}), whose entries depend polynomially on finitely many input coefficients (Theorem \texttt{thm:ancestry}). That is depth $0$, where the outer-block construction of this article, which needs $n\ge1$, does not apply; particular solutions of forced equations are not treated there.
% ed. (2026-09-29): on filing batch 50, "The question as posed, at depth n>=1 with the new logarithm T and with forcing, remains open." was replaced: a later package answers it.
At depth $n\ge1$, with the new logarithm $T$ and with forcing, the question is answered for this article's whole class by the later package \path{Analysis/Transseries/docs/series-and-transseries/Exact_Logarithmic_Degree_Smith_Invariants/} (batch~50; not reviewed): the Smith exponents of a lifted residue series satisfy $\sum\nu_i=r$, $0\le\nu_i\le s$, give the whole filtration, and are decided by block Toeplitz ranks of a finite jet, with rational obstruction certificates; the Jordan-shift family of \cref{thm:sharp-matrix} is its single-block case.
\end{ednote}

\subsection{Cancellation-sensitive resonance graphs}
The semigroup condition \eqref{eq:selection} detects impossible matrix entries but does not prove that a permitted path contributes nontrivially. Different derivative orders and intermediate blocks may cancel. Can a weighted directed graph of residue transfers characterize the leading nonvanishing path and sharpen both the cutoff and the logarithmic-degree bound?

For monomial coefficients one can first classify which edge weights depend only on $P'(\lambda)$ and which require higher residue moments. A useful theorem would distinguish structural zeros from coefficient-dependent cancellation without expanding all lower logarithmic terms.

% ed. (2026-09-29): added on filing batch 50; a partial answer for a scalar subclass.
\begin{ednote}\label{ed:resonance-graph}
For operators $P(E)+\sigma\sum_\eta x^{-\eta}q_\eta(E)$ with simple real roots, the later package \path{Analysis/Transseries/docs/series-and-transseries/Exact_Logarithmic_Degree_Smith_Invariants/} (batch~50; not reviewed) shows that the relevant weighted graph is the matrix $D_P^{-1}N$, $N_{ij}=q_{\lambda_j-\lambda_i}(\lambda_j)$, whose Jordan block sizes are the degrees; a fourth-order family has the same unweighted graph and different degrees. The general case and the sharpening of the cutoff remain open.
\end{ednote}

\subsection{Perturbations small only in logarithms}
The hypothesis $\dprim(a_j)\le-\eps$ excludes coefficients such as $1/\log x$, whose outer exponent is zero. In that setting a perturbation can stay in one outer root block indefinitely. Can a multi-level version of the reduction, contracting successively in $x$, $\log x$, and deeper logarithms, still yield finite obstruction data?

The present proof cannot simply replace ``outer-small'' by ``small'' in the full lexicographic order. Its finite cutoff and its rule that a root is visited at most once would no longer apply. A precise next step is a classification of single-root operators with coefficients supported in negative powers of $L$, separating those needing only $T$ from those requiring exponentials or further logarithmic depth.

% ed. (2026-09-29): added on filing batch 53; the first-order case of this classification, and operators that factor.
\begin{ednote}\label{ed:log-small}
For first-order operators the classification asked for is given by the later package \path{Analysis/Transseries/docs/series-and-transseries/Triangular_Hahn_Differential_Systems_Without_Smallness/} (batch~53; not reviewed; \texttt{triangular\_hahn\_systems.tex}, Theorems \texttt{thm:rankone} and \texttt{thm:descent}). With $L=\log x$, a coefficient $x^{-1}\bigl(\lambda+\sum_{\beta>0}c_\beta L^{-\beta}\bigr)$ is a logarithmic derivative in $\HH_n$ exactly when no exponent $\beta$ with $c_\beta\ne0$ lies in $(0,1)$, so always for integer powers; then only $T$ is needed, with degree at most one. Otherwise, as for $x^{-1}L^{-\alpha}$ with $0<\alpha<1$, the operator is bijective on $\HH_n$ and has no nonzero homogeneous solution at any finite logarithmic depth; the actual solution $\exp(L^{1-\alpha}/(1-\alpha))$ is an exponential of a large series. Further logarithmic depth never creates a homogeneous solution. Higher-order single-root operators are covered by its triangular theorem (\texttt{thm:triangular}, \texttt{cor:gauge}) only when they factor into first-order operators over the field; the rest of the question remains open.
\end{ednote}

\subsection{Parameter-dependent root collisions}
When two distinct real roots coalesce, $s$ decreases but a multiplicity increases. The degree theorem predicts less growth in the new logarithm, while the natural coefficient formulas may become singular in the root difference. Can one construct a parameter-uniform block basis and obstruction matrix that pass continuously, or in a controlled confluence sense, through such a collision?

A test case is $P_\tau(z)=z(z-\tau)$ with a monomial outer-small perturbation. The fixed-parameter theorem is already proved here; the unresolved target is a uniform representation and an interpretation of the divergent coefficients as $\tau\to0$.

\subsection{Matrix systems and nonreal indicial roots}
For a matrix polynomial $P(E)$, kernels at a root may contain Jordan chains before any perturbation. A matrix-valued residue complex should replace the scalar moments. Which invariant controls the new-logarithm degree: the number of distinct root levels, the lengths of internal Jordan chains, or a combination?

Nonreal scalar roots raise a different issue. The present real monomial field deliberately omits $x^{i\beta}$ and oscillation. An extension with complex powers would need a careful support convention rather than an artificial total ordering of oscillatory functions. A first goal is an algebraic version for a fixed finite collection of imaginary exponents, without claiming an ordered Hardy-field interpretation.

% ed. (2026-09-29): records the later package that answers the first
% paragraph for real-spectrum first-order systems at depth 0.
\begin{ednote}\label{ed:matrix-systems}
% ed. (2026-09-29): "the previous note" became ambiguous once other notes were inserted between; the note is now named.
For first-order systems at depth $0$ with real spectrum (the setting of the note after ``The exact minimal power of the promoted logarithm''), the later package \path{Analysis/Transseries/docs/series-and-transseries/Hahn_Fuchsian_Resonance_Analytic_Normalization/} identifies the invariant: the new-logarithm degree of a fundamental matrix is exactly $\nu(B)-1\le m-1$, where $\nu(B)$ is the nilpotency index of $B=N+\sum_\rho R_\rho$, which combines the internal Jordan chains (through $N$) with the passage between root levels (through the resonant blocks $R_\rho$) (\texttt{hahn\_fuchsian.tex}, Theorem \texttt{thm:logs}; the bound is attained in Example \texttt{ex:chain}, the analogue of the Jordan-shift family of \cref{thm:sharp-matrix}). When $N=0$, $B$ is strictly block-triangular over the $s$ distinct eigenvalue levels, so the degree is at most $s-1$, the homogeneous bound of \cref{thm:promotion}; that package does not state this consequence. Matrix polynomials $P(E)$ of higher order, depth $n\ge1$, and nonreal roots are not treated there.
% ed. (2026-09-29): sentence added on filing batch 53.
For triangular first-order systems at every depth, with arbitrary coefficients, the later package \path{Analysis/Transseries/docs/series-and-transseries/Triangular_Hahn_Differential_Systems_Without_Smallness/} (batch~53; not reviewed) bounds the new-logarithm degree by the number of split factors on a coupling path and classifies it by the Smith exponents of a residue series with determinant $z^r$; nontriangular systems and nonreal roots are not treated there either.
\end{ednote}

\subsection{Nonlinear equations and an effective obstruction variety}
Consider $P(E)u+Ru+F(u)=f$, with a nonlinear term that is suitably small in a support filtration. Eliminating the infinite-dimensional complement may leave a finite nonlinear equation in the kernel coordinates. Can the linear matrix $M$ be replaced by a convergent formal obstruction map, with explicit support bounds for its Taylor coefficients?

The existing generalized-series implicit-function theory is a natural framework \cite{vdh2001}. The specific next theorem would retain resonance coordinates rather than assume an everywhere invertible linearization, and would state which nonlinearities preserve a finite logarithmic-depth bound.

\subsection{Effective coefficient representations}
For arbitrary real Hahn coefficients, coefficient extraction and equality need not be computable. What practical representation makes the finite-depth cutoff of \cref{thm:cutoff} an actual terminating symbolic algorithm?

A useful restricted target combines rational exponent lattices, computable well-based supports, and coefficient blocks generated by a finite differential-algebraic grammar. One should specify the exact oracle assumptions and prove complexity in the number of reachable root differences, rather than describe the finite matrix as automatically computable.

\subsection{Analytic realization beyond the displayed block}
The real construction in \cref{sec:analytic} certifies the decisive outer-zero block, not every subsequent block. Can one show that the same Volterra solutions have the entire formal Hahn expansion, with explicit truncation bounds simultaneously in powers of $x^{-1}$ and $(\log x)^{-1}$?

For a broader class, impose analytic or Gevrey bounds on the coefficient blocks and ask whether the finite residue matrix identifies precisely which logarithmic terms occur in sectorial actual solutions. A particularly useful result would translate \eqref{eq:tail} into an analytic residual bound without silently treating formal support order as numerical size.

\subsection{Other ordered exponent groups}
Our outer exponents are real, and the uniform $\eps$ gives a simple Archimedean cutoff. A more general ordered group may have positive infinitesimals or no useful integer-valued depth bound. Which version of ``strict decrease'' guarantees summability of the geometric operator inverse and finiteness of its projection to the resonant part?

One should separate two demands: existence of the full formal inverse and termination of the finite-dimensional obstruction calculation. They coincide under the present real uniform-drop hypothesis, but need not coincide in a non-Archimedean exponent group.

% ed. (2026-09-29): cross-link to a related, not answering, later package.
\begin{ednote}\label{ed:ordered-groups}
The later package \path{Analysis/Transseries/docs/series-and-transseries/Reversion_Beyond_Archimedean_Valuations/} draws the analogous separation for compositional reversion over an arbitrary ordered exponent group: strong Hahn convergence always holds, uniform valuation convergence holds exactly when the multiples of the least support exponent are cofinal, and every finite coefficient request factors through a finite-rank quotient. It concerns reversion of near-identity maps, not the differential operators of this article, so this question remains open.
\end{ednote}

\subsection{A reusable proof-assistant library for finite defects}
The identities $P_0G=\id-J\Pi$ and $GP_0=\id-K\Lambda$ package an infinite-dimensional problem into a finite kernel/cokernel interface. Can a library expose this interface independently of the Hahn implementation and reuse it for difference equations, summation, or sequence inversion?

The immediate target is a formally checked version of \cref{thm:fredholm} with explicit summability hypotheses and then one concrete Hahn instance. This would be a meaningful extension of the repository's existing block-antiderivative results, rather than a new count of unrelated declarations.

\section{Conclusion}
At finite logarithmic depth, the obstruction to solving a differential equation need not be an uncontrolled infinite tail. For outer-small perturbations of $P(xD)$, it is a finite residue matrix, and only finitely many perturbation depths can affect that matrix. Its rank measures the exact discrepancy between solvability at the original depth and solvability after promotion.

One more logarithm always suffices in this class, but the degree in that logarithm records the chain of distinct resonant outer exponents. The Jordan-shift family proves that arbitrarily high powers can be forced while only one new logarithmic level is needed. The second-order member demonstrates that the formal obstruction can also be accompanied by actual real solutions and a quantitative bound.

The outcome is a specific extension of the repository's polynomial-block calculus: a concrete Hahn ambient, explicit residue moments, finite-dimensional obstruction data, sharp depth/degree control, and a separate analytic certificate. The broader nonlinear, parameter-uniform, and general analytic-realization questions remain distinct research tasks.

\appendix
\section{A compact abstract form of the reduction}
\label{app:abstract}
For reuse outside transseries, let $X$ be a real vector space, $V$ finite-dimensional, and suppose linear maps
\[
 P_0,G:X\to X,\quad K,J:V\to X,\quad\Lambda,\Pi:X\to V
\]
satisfy
\[
 P_0K=0,\quad\Pi P_0=0,\quad\Lambda K=\Pi J=\id,
\]
\[
 P_0G=\id-J\Pi,\quad GP_0=\id-K\Lambda,
 \quad\Lambda G=GJ=0.
\]
For any $R:X\to X$ such that $\id+GR$ is invertible, define $\shift$, $M$, and $b$ as in \eqref{eq:neumann} and \eqref{eq:effective-data}, interpreting $\shift$ simply as the inverse if no infinite sum is involved. The proof of \cref{thm:fredholm} gives
\[
 \ker(P_0+R)\cong\ker M,\qquad
 \coker(P_0+R)\cong\coker M.
\]
No norm, topology, or differential structure is needed for this algebraic statement. The Hahn-specific work is to construct the splitting maps, establish the geometric inverse, and show that its finite-dimensional projection terminates. This decomposition is also why the finite-dimensional exact tests in the supplied script are relevant but cannot certify the Hahn support arguments.

\section{Support facts used in the constructions}
\label{app:support}
We spell out the elementary local-finiteness principle used repeatedly. Suppose $(S_q)_{q\ge0}$ are reverse well-ordered subsets of $\R^{n+1}$ in lexicographic order and every first coordinate in $S_q$ is at most $A-q\eps$, where $\eps>0$. Then $S=\bigcup_qS_q$ is reverse well-ordered. To see this, choose one element $a$ of a nonempty subset of $S$. Only finitely many $S_q$ can contain an element with first coordinate at least $a_0$. Each of their intersections with the subset has a greatest element, and the greatest of those finitely many choices is greatest in the whole subset. Every fixed exponent lies in only finitely many $S_q$ as well.

If each $S_q$ is itself the support union of a summable family, the finite-contribution condition at a fixed exponent then proves summability of the doubly indexed family. The same argument applies inside an outer block with $L$ in place of $x$ and $\eps=1$. It proves the rearrangement assertions for all geometric and formal functional-calculus sums used in the paper.

For ordinary Hahn multiplication, finite sums of reverse well-ordered supports are reverse well-ordered, and each resulting exponent has finitely many decompositions. One way to check the obstruction to an infinite bad sequence is to use the following elementary fact: a sequence in a reverse well-ordered total order has an infinite nonincreasing subsequence. If infinitely many indices are at least as large as all their successors, choose those; otherwise construct a forbidden infinite increasing subsequence. Given an increasing sequence of sums, first pass to a nonincreasing subsequence in one summand support and then in the other, obtaining a contradiction. Given infinitely many distinct decompositions of one fixed exponent, a strictly decreasing subsequence in one coordinate would force a strictly increasing one in the other. This proves the required local finiteness. Finite unions are immediate.

These arguments explain why none of the infinite operator expressions in the article is being justified merely by suggestive notation or by an unproved interchange of infinite sums.

\section{Package contents and reproducibility}
The delivered archive contains the editable source \texttt{residue\_fredholm.tex}, the compiled article \texttt{residue\_fredholm.pdf}, a build/readme file, the verification script, and its recorded JSON output. The document uses standard \LaTeX{} packages and can be built with
\begin{verbatim}
pdflatex -interaction=nonstopmode -halt-on-error residue_fredholm.tex
pdflatex -interaction=nonstopmode -halt-on-error residue_fredholm.tex
pdflatex -interaction=nonstopmode -halt-on-error residue_fredholm.tex
python verification/verify.py
\end{verbatim}
The verification script uses rational arithmetic for symbolic checks and requires \texttt{sympy} and \texttt{mpmath}. The exact output counts and numerical illustration values described in \cref{sec:verification} refer to the run included with this package. They do not assert a proof-assistant verification of the mathematical manuscript.

\begin{thebibliography}{10}
\small
\bibitem{repo-root}
Vladimir Reshetnikov, \emph{ProveIt: Transseries project README}, repository snapshot
\texttt{8813fb3a2cdd1702d3ea7f7ca5d63ae710630bef}, inspected 29 September 2026.
\url{https://github.com/VladimirReshetnikov/ProveIt/blob/8813fb3a2cdd1702d3ea7f7ca5d63ae710630bef/Analysis/Transseries/README.md}.

\bibitem{repo-volume}
Vladimir Reshetnikov, \emph{ProveIt: Series and transseries documentation and canonical-volume status}, same snapshot and access date.
\url{https://github.com/VladimirReshetnikov/ProveIt/blob/8813fb3a2cdd1702d3ea7f7ca5d63ae710630bef/Analysis/Transseries/docs/series-and-transseries/README.md}.
See also the adjacent \texttt{Transseries\_And\_Inversion/README.md} for the declaration-level crosswalk and explicit scope boundaries.

\bibitem{repo-block}
Vladimir Reshetnikov, \emph{ProveIt: TransseriesBlockAntiderivative.lean}, same snapshot and access date. Inspected file blob:
\texttt{1a4b5e335b5429a267bcfa4c0c9b35aadf9be527}.
\url{https://github.com/VladimirReshetnikov/ProveIt/blob/8813fb3a2cdd1702d3ea7f7ca5d63ae710630bef/Analysis/Transseries/Lean/Transseries/TransseriesBlockAntiderivative.lean}.

\bibitem{km2012}
Salma Kuhlmann and Micka\"el Matusinski,
\emph{Hardy type derivations on generalised series fields},
Journal of Algebra \textbf{351} (2012), 185--203.
DOI: \href{https://doi.org/10.1016/j.jalgebra.2011.11.024}{10.1016/j.jalgebra.2011.11.024}.
Preprint: \href{https://arxiv.org/abs/0903.2197}{arXiv:0903.2197}.

\bibitem{vdh2001}
Joris van der Hoeven, \emph{Operators on generalized power series},
Illinois Journal of Mathematics \textbf{45} (2001), 1161--1190.
DOI: \href{https://doi.org/10.1215/ijm/1258138061}{10.1215/ijm/1258138061}.
Author's text: \url{https://www.texmacs.org/joris/noeth/noeth.html}.

\bibitem{vdh2006}
Joris van der Hoeven, \emph{Transseries and Real Differential Algebra},
Lecture Notes in Mathematics \textbf{1888}, Springer, 2006.
DOI: \href{https://doi.org/10.1007/3-540-35590-1}{10.1007/3-540-35590-1}.

\bibitem{hogancamp}
Matthew Hogancamp, \emph{Homological perturbation theory with curvature},
preprint, 2019.
\href{https://arxiv.org/abs/1912.03843}{arXiv:1912.03843}.
Cited for the broader perturbative framework, not as a source for the particular Hahn residue and degree formulas proved here.
\end{thebibliography}
\end{document}
